\documentclass[11pt]{article}
\usepackage[margin=1in]{geometry}
\usepackage{amsmath,amssymb}
\usepackage{bm}
\usepackage{hyperref}

\title{Method Description for the MFA-NB xVERSE Model}
\author{}
\date{}

\begin{document}
\maketitle

\section{Motivation}

The goal of the model is to learn a latent representation that is simultaneously useful for biological embedding, robust to incomplete gene panels, and capable of generating realistic single-cell expression profiles. A standard variational autoencoder with an isotropic Gaussian prior provides a smooth latent space and a well-defined generative model, but its prior geometry is often too simple for whole-atlas single-cell data. Across tissues, diseases, and cellular states, the latent distribution is expected to contain multiple major modes as well as structured local variation around each mode. We therefore use a mixture-of-factor-analyzers (MFA) prior to give the latent space an explicit geometric structure.

The desired embedding has several properties. First, cells with similar biological identity should remain nearby even when observed through different gene panels or random masks. Second, the representation should preserve rare or fine-grained biological variation rather than only capturing dominant high-expression programs. Third, the latent space should be generative: sampling from the prior should produce a valid expression profile through the decoder. Finally, the latent geometry should be interpretable: mixture components can be viewed as coarse cellular prototypes, while low-rank factor directions within each component can capture continuous intra-state variation.

\section{MFA Prior}

Let $z \in \mathbb{R}^d$ denote the latent variable used by the decoder. The prior is a mixture of factor analyzers. For each cell, the generative process first samples a discrete component index
\begin{equation}
    c \sim \mathrm{Cat}(\pi),
\end{equation}
where $\pi$ is the learned mixture weight vector. Conditional on component $c=k$, the model samples a low-dimensional factor
\begin{equation}
    u \sim \mathcal{N}(0, I_R),
\end{equation}
and a diagonal residual noise term
\begin{equation}
    \epsilon \sim \mathcal{N}\left(0, \mathrm{diag}\left(\exp(\ell_k)\right)\right),
\end{equation}
where $\ell_k$ is the learned log-variance vector for component $k$. The latent variable is then generated as
\begin{equation}
    z = \mu_k + A_k u + \epsilon.
\end{equation}
Here $\mu_k \in \mathbb{R}^d$ is the component center and $A_k \in \mathbb{R}^{d \times R}$ is the component-specific factor loading matrix. Marginally,
\begin{equation}
    p(z \mid c=k) = \mathcal{N}\left(\mu_k, A_k A_k^\top + \mathrm{diag}\left(\exp(\ell_k)\right)\right),
\end{equation}
and therefore
\begin{equation}
    p(z) = \sum_{k=1}^{K} \pi_k \, \mathcal{N}\left(z; \mu_k, A_k A_k^\top + \mathrm{diag}\left(\exp(\ell_k)\right)\right).
\end{equation}

This prior is more expressive than a single Gaussian but remains structured. The component centers define coarse regions of the latent landscape. The low-rank matrices $A_k$ define dominant local directions of variation around each prototype. The diagonal residual term captures remaining component-specific noise. This structure is appropriate for single-cell data because cellular states are neither purely discrete nor fully unstructured: major cell identities form broad modes, while activation, differentiation, disease response, and technical effects often vary continuously within each mode.

\section{Encoder and Approximate Posterior}

The encoder receives two pieces of information: the observed expression vector and the gene-observation mask. Counts are transformed by $\log(1+x)$ before entering the encoder, and the mask indicates which genes are observed for the cell. This is important because different assays or panels observe different subsets of genes; the model must distinguish true zero expression from missing genes.

The encoder produces a hidden representation $h(x,m)$, where $x$ is the count vector and $m$ is the gene mask. The approximate posterior follows the MFA structure:
\begin{equation}
    q(c,u,\epsilon \mid x,m)
    = q(c \mid x,m)\, q(u \mid x,m,c)\, q(\epsilon \mid x,m,c).
\end{equation}
The component posterior is categorical,
\begin{equation}
    q(c \mid x,m) = \mathrm{Cat}\left(\rho(x,m)\right),
\end{equation}
where $\rho$ is predicted by a neural posterior head. For each component $k$, the posterior over the MFA factor and residual variable is diagonal Gaussian:
\begin{align}
    q(u \mid x,m,c=k) &= \mathcal{N}\left(m^u_k(x,m), \mathrm{diag}\left(\exp(s^u_k(x,m))\right)\right), \\
    q(\epsilon \mid x,m,c=k) &= \mathcal{N}\left(m^\epsilon_k(x,m), \mathrm{diag}\left(\exp(s^\epsilon_k(x,m))\right)\right).
\end{align}
The resulting component-specific posterior over $z$ is induced by
\begin{equation}
    z = \mu_k + A_k u + \epsilon.
\end{equation}
Thus, the full posterior over $z$ is a mixture over components. The model also computes a moment-matched posterior mean for diagnostics and embedding extraction:
\begin{equation}
    \mathbb{E}_q[z \mid x,m]
    = \sum_{k=1}^{K} q(c=k \mid x,m)
    \left(\mu_k + A_k m^u_k(x,m) + m^\epsilon_k(x,m)\right).
\end{equation}

\section{Decoder and Negative Binomial Likelihood}

The model is generative: expression profiles are generated from the latent variable $z$. The decoder maps $z$ to gene-level logits and a library-size term. The gene probabilities are obtained by a softmax over genes,
\begin{equation}
    p_g(z) = \frac{\exp(a_g(z))}{\sum_{g'} \exp(a_{g'}(z))},
\end{equation}
where $a_g(z)$ is the decoder logit for gene $g$. The library-size head predicts a positive total count scale $L(z)$, and the mean expression rate for gene $g$ is
\begin{equation}
    r_g(z) = L(z) p_g(z).
\end{equation}

Counts are modeled with a negative binomial likelihood:
\begin{equation}
    x_g \mid z \sim \mathrm{NB}\left(r_g(z), \theta_g\right),
\end{equation}
where $r_g(z)$ is the mean and $\theta_g$ is a learned dispersion parameter. The negative binomial likelihood is used because single-cell count data are overdispersed relative to a Poisson model. This allows the decoder to model both the expected expression profile and gene-specific count variability.

When a cell is observed through an incomplete panel, reconstruction loss is evaluated only on observed genes:
\begin{equation}
    \log p(x_{\mathrm{obs}} \mid z) = \sum_{g:m_g=1} \log \mathrm{NB}\left(x_g; r_g(z), \theta_g\right).
\end{equation}
This avoids penalizing the model for genes that were not measured in a given panel.

\section{Training Objective}

The variational objective contains a reconstruction term and KL terms for the structured posterior:
\begin{align}
    \mathcal{L}_{\mathrm{ELBO}}
    =&\; \mathbb{E}_{q(c,u,\epsilon \mid x,m)}\left[\log p(x_{\mathrm{obs}} \mid z)\right] \\
    &- \mathrm{KL}\left(q(c \mid x,m) \;\|\; p(c)\right) \\
    &- \mathbb{E}_{q(c \mid x,m)}
       \mathrm{KL}\left(q(u \mid x,m,c) \;\|\; p(u)\right) \\
    &- \mathbb{E}_{q(c \mid x,m)}
       \mathrm{KL}\left(q(\epsilon \mid x,m,c) \;\|\; p(\epsilon \mid c)\right).
\end{align}
In minimization form, the core loss is
\begin{equation}
    \mathcal{J}_{\mathrm{VAE}}
    = \mathcal{J}_{\mathrm{recon}}
    + \beta_{\mathrm{KL}}\left(
        \mathcal{J}_{c}
        + \alpha_u \mathcal{J}_{u}
        + \alpha_{\epsilon} \mathcal{J}_{\epsilon}
      \right),
\end{equation}
where $\beta_{\mathrm{KL}}$ controls the strength of the prior constraint and $\alpha_u, \alpha_{\epsilon}$ optionally reweight the factor and residual KL terms. In practice, this objective can be augmented with auxiliary terms for mask invariance and label-guided alignment, but these are auxiliary regularizers rather than part of the pure likelihood model.

\section{Mask Invariance and Panel Robustness}

To make the representation robust to gene-panel differences, the model is trained with random mask augmentation. Two masked views of the same cell can be generated from the same original observed gene set. The reconstruction objective asks each partial view to recover the observed panel genes. In addition, a contrastive consistency term can align the posterior distributions inferred from different masked views of the same cell. This encourages the encoder to represent cell identity rather than the accidental choice of observed genes.

The distribution-level contrast compares posterior summaries $q(z \mid x,m_1)$ and $q(z \mid x,m_2)$ for two masks $m_1$ and $m_2$ of the same cell. Rather than forcing two stochastic samples of $z$ to be identical, the contrastive objective encourages the inferred posterior distributions to agree. This is useful because the model should tolerate uncertainty from partial observations while maintaining a stable biological representation.

\section{Model Inputs and Design Considerations}

The model uses raw count vectors, gene-observation masks, sample or batch identifiers when batch conditioning is enabled, and optional cell-type annotations for auxiliary supervision. The count vector provides the expression signal; the mask tells the model which genes are measured; batch information can be used by the decoder to account for batch-specific count structure when desired. The core latent representation itself is learned from expression and mask information through the encoder.

The design reflects four main considerations. First, a geometrically structured prior is needed to describe the multimodal and continuous nature of whole-atlas single-cell data. Second, the embedding must be robust to panel missingness and random gene masking. Third, the model must remain generative, so the latent variable $z$ is decoded into a full negative-binomial expression distribution. Fourth, the prior should be interpretable: mixture centers represent coarse prototypes, factor directions represent local axes of variation, and sampling from these regions produces expression profiles through the same decoder used during training.

\end{document}
